\section{Reproducing this document (toolchain)}\label{sec:repro}

This whitepaper is itself a small, self-contained build: a contributor who clones the
repository should be able to regenerate the PDF, with every diagram, from the sources
checked into \srcf{docs/PartPlacementDesignLatex/}. This appendix documents that build so
that the figures referenced throughout the preceding sections remain reproducible rather
than opaque binary artifacts. The guiding principle is that nothing rendered into the final
document is committed in its rendered form: the figures are derived, on demand, from text
sources by a deterministic pipeline, so a reviewer can audit \emph{how} every diagram was
produced and regenerate it.

\subsection{Directory layout}

The document is assembled from a small number of well-separated source kinds. The master
file \srcf{PartPlacement.tex} carries only the document scaffolding---title block, front
matter, and an ordered list of \cppinline|\input| directives---while the prose lives in
\srcf{sections/}, one file per section (this appendix is
\srcf{sections/A1-reproducibility.tex}). The shared style contract is isolated in
\srcf{preamble.tex}, which defines the colour palette, the fonts (XCharter text,
\code{newtxmath}, Inconsolata mono), every callout environment, and the domain macros
(\code{\textbackslash src}, \code{\textbackslash classfig}, \code{\textbackslash tikzfig},
and the rest) that the section files depend on. Keeping the preamble separate is what lets
each section file be authored and reviewed in isolation: a section is just a body, never a
standalone \LaTeX{} program.

\begin{designrule}[Two diagram toolchains, by figure kind]
Class diagrams are authored in PlantUML and live as \srcf{diagrams/*.puml}; activity,
component, geometry, and worked-example figures are hand-built in Ti\emph{k}Z and live as
\srcf{diagrams/tikz/*.tex}. The split is deliberate: PlantUML excels at laying out
box-and-line class structure with no manual coordinate work, whereas the flow, axis, and
geometry figures need pixel-level control that is far easier to express directly in
Ti\emph{k}Z than to coerce out of a layout engine.
\end{designrule}

The remaining directories are mechanical. \srcf{diagrams/out/} holds the rendered diagram
artifacts (the intermediate SVGs and the PDFs that \LaTeX{} actually includes); it is a
build output and is reconstructed from sources. \srcf{tools/plantuml.jar} is the pinned
PlantUML renderer, fetched on demand rather than vendored as a hand-managed binary. The
shared PlantUML theme \src{diagrams/\_style.iuml}{1} is given the \code{.iuml} extension
precisely so that it stays out of the Makefile's \srcf{diagrams/*.puml} glob---it is an
include, not a diagram, and must not be rendered as one.

\subsection{The rendering pipeline}

Diagrams flow through three stages before they reach the document. Each PlantUML source is
first rendered to SVG, the SVG is converted to a vector PDF, and that PDF is pulled into the
document by \code{\textbackslash includegraphics} (via the \code{\textbackslash classfig}
macro). Schematically:

\begin{minted}{text}
diagrams/*.puml      --(plantuml -tsvg)-->     diagrams/out/*.svg
diagrams/out/*.svg   --(rsvg-convert -f pdf)-->  diagrams/out/*.pdf
diagrams/out/*.pdf + sections/*.tex  --(latexmk -pdf)-->  PartPlacement.pdf
\end{minted}

SVG is used as the intermediate because it is PlantUML's cleanest vector output; the
SVG-to-PDF step uses \code{rsvg-convert} so the figures embed as resolution-independent
vectors rather than rasterized bitmaps. The converter is configurable---the Makefile exposes
a \code{SVG2PDF} variable (\src{Makefile}{33}) so a contributor without \code{rsvg-convert}
can substitute, for example, an Inkscape export command without editing the rules. The
Ti\emph{k}Z figures need no such pipeline: they are \LaTeX{} sources included directly by
\code{\textbackslash tikzfig} and compiled in line by the main engine.

The document engine is \code{pdflatex}, driven by \code{latexmk} so that the multi-pass
dance (cross-references, the table of contents, \code{cleveref}, and the bookmark tree)
converges automatically. The one non-default requirement is \code{-shell-escape}: the
listings throughout this document are typeset by \code{minted}, which shells out to the
Python \code{pygmentize} tool to perform syntax highlighting. As the preamble notes at
\src{preamble.tex}{3}, \code{minted} cannot run without shell-escape, so the Makefile's
\code{latexmk} invocation enables it unconditionally (\src{Makefile}{35}).

\begin{keyinsight}[PlantUML resolves \texttt{!include} relative to the working directory]
Every \srcf{*.puml} begins with \cppinline|!include _style.iuml| to pull in the shared theme
(\src{diagrams/class-current.puml}{4}). PlantUML resolves that path relative to the
\emph{current working directory}, not the file being rendered, so running it from the
repository root would fail to find the theme. The Makefile therefore \code{cd}s into
\srcf{diagrams/} before invoking the renderer, passing the jar and the output directory as
absolute paths (\src{Makefile}{60}). This is the single most common foot-gun when rebuilding
a diagram by hand: invoke PlantUML from \srcf{diagrams/}, or the \code{!include} will not
resolve.
\end{keyinsight}

\subsection{Authoring format: a sample class diagram}

A PlantUML class diagram is a compact text description of types and their relations; the
renderer handles all layout. \Cref{lst:puml-sample} shows the head of
\srcf{diagrams/class-current.puml} (the source behind \cref{fig:class-current}): it includes
the shared theme, sets a title, and declares the abstract \code{Part} with its members and
the stereotype that drives its colour. The full file goes on to declare the concrete part
subclasses, the composition edge, and an explanatory note, but the excerpt is enough to show
that the diagram is fully auditable text---there is no hand-placed geometry to drift out of
sync with the code it depicts.

\begin{listing}[htbp]
\begin{minted}{text}
@startuml class-current
!include _style.iuml
title Current model::part composite tree — CM-to-CM child offset

abstract class "Part" as Part <<abstract>> {
  +using Id = std::uint64_t
  --
  -cm : Vector3
  -needsRecomputing : bool
  -childParts : vector<tuple<shared_ptr<Part>, Vector3>>
  --
  +getId() : Id
  +addChildPart(child, Vector3 position)
  +getChildParts() : const vector<tuple<..>>&
  #computeCompositeAt(t)
}
Part "1" o-- "childParts *" Part : composite (by move)
@enduml
\end{minted}
\caption{Head of \texttt{diagrams/class-current.puml}: a PlantUML class diagram is plain
text, themed by \texttt{!include \_style.iuml} and rendered to SVG then PDF by the Makefile.}
\label{lst:puml-sample}
\end{listing}

The shared theme \src{diagrams/\_style.iuml}{21} maps stereotypes (\code{<<abstract>>},
\code{<<concrete>>}, \code{<<pod>>}, \code{<<newtype>>}, \code{<<enumc>>}, \code{<<freefn>>},
\code{<<legacy>>}) onto the same colour palette as the body text, so the class diagrams and
the prose share one visual language. Per its own header comment, this PlantUML build rejects
the block form \code{skinparam class<<stereo>> \{ \}} (it exits with status 200); the theme
therefore uses the supported \code{Property<<stereo>>} suffix form inside a single
\code{skinparam class \{ \}} block (\src{diagrams/\_style.iuml}{5}). A contributor adapting
the theme should preserve that form.

\subsection{Commands}

The Makefile wires the whole pipeline into a handful of phony targets. From
\srcf{docs/PartPlacementDesignLatex/}:

\begin{minted}{text}
make            # default: render all diagrams, then build PartPlacement.pdf
make diagrams   # render only the PlantUML class diagrams to diagrams/out/*.pdf
make tools      # fetch tools/plantuml.jar (pinned version) if it is missing
make clean      # remove LaTeX aux files, the minted cache, and diagram SVGs
make veryclean  # clean, and also remove PartPlacement.pdf + rendered diagram PDFs
\end{minted}

The dependency graph is arranged so that a clean checkout builds in a single \code{make}
invocation: the PDF target depends on the rendered diagram PDFs (\src{Makefile}{67}), each
diagram PDF depends on its SVG, and each SVG depends on its \srcf{*.puml} source, the shared
theme, and---as an order-only prerequisite---the PlantUML jar. The jar rule
(\src{Makefile}{50}) fetches a pinned PlantUML version over \code{curl} into \srcf{tools/}
only when the jar is absent, so the network is touched at most once and the version is
reproducible. The intermediate SVGs are marked \code{.SECONDARY} (\src{Makefile}{42}) so
\code{make} does not delete them as transient by-products between the two diagram stages.

\code{make clean} runs \code{latexmk -c} to drop the \LaTeX{} auxiliary files, then removes
the \code{minted} cache (the \srcf{\_minted-*} directories that hold the cached highlighter
output) and the intermediate SVGs, while leaving the finished PDF and the rendered diagram
PDFs in place for inspection; \code{make veryclean} additionally runs \code{latexmk -C} and
deletes the final PDF and the entire \srcf{diagrams/out/} tree, returning the directory to a
pristine, source-only state.

\subsection{Prerequisites}

To rebuild from a clean checkout a contributor needs: a Java runtime (\code{java}, to run
the fetched \srcf{tools/plantuml.jar}); \code{curl} (to fetch the jar on first build);
\code{rsvg-convert} or a substitute set via \code{SVG2PDF} (for the SVG-to-PDF step); and a
\TeX{} distribution providing \code{latexmk}, \code{pdflatex}, and the \code{minted} package
together with its \code{pygmentize} helper. With those present, \code{make} reproduces this
document end to end.
